\section{A segunda camada de evidência: hierarquia e eficiência}
\label{sec:fase2}

A segunda fase da estimação atacou a mesma pergunta com três famílias
de métodos ainda não usadas no documento: modelos
\textbf{hierárquicos} (que tratam os 11 anos de cada hospital como
parentes, não como observações soltas), a análise envoltória de dados
(\textbf{DEA}, que constrói uma fronteira de melhores práticas e mede
a distância de cada hospital até ela) e a fronteira estocástica
(\textbf{SFA}, a contraparte paramétrica do DEA). O pacote completo de
resultados, com scripts, diagnósticos e um relatório de execução
auditável, está na pasta resultados\_fase2 do repositório; aqui ficam
os achados que mudam ou reforçam a leitura do documento.

\textbf{O que a camada hierárquica confirma.} Estimados com um
efeito próprio para cada hospital, o tempo de permanência e o faturamento
real reproduzem quase exatamente os resultados da
Seção~\ref{sec:resultados} (queda de 10,2\% no TMP; faturamento sem
diferença), a produção continua maior nas OSS (33\%, com intervalo
inteiro acima do zero) e a ocupação aponta na mesma direção com
magnitude menor (11,5\%), \textbf{mas com intervalo que cruza o zero na
especificação hierárquica}, em ambos os casos porque parte do
contraste entre as redes é absorvida pelo efeito de hospital; e a mortalidade aparece
\textbf{0,262 menor no logito} (intervalo de 0,410 a 0,114 negativos;
na escala da própria taxa, isso equivale a cerca de \textbf{1,3
ponto percentual a menos de mortalidade}, na altura da mediana do
painel, que é a forma interpretável do mesmo número),
com precisão maior que a da
Seção~\ref{sec:resultados}.\footnote{Duas especificações alternativas
foram estimadas para esse mesmo coeficiente e não são reportadas aqui
por apontarem na mesma direção e magnitude: sem o termo de Mundlak,
0,233 negativo; com efeito aleatório de ano, 0,260 negativo. A tabela
completa está em \texttt{tabB\_sensibilidade\_mortalidade.csv}.} O teste
decisivo, porém, é outro: \textbf{retirando os cinco conversores, a
diferença de mortalidade entre as redes desaparece}: o coeficiente vai
a 0,003 positivo, com intervalo de 0,414 negativo a 0,420 positivo,
cruzando o zero com folga, enquanto TMP, produção, ocupação e
eficiência não só sobrevivem como se \textbf{fortalecem}: sem os
conversores o TMP vai de 0,108 para 0,281 negativo, a produção de 0,282
para 0,492, a ocupação de 0,108 para 0,218 (passando a excluir o zero,
o que não fazia no modelo principal) e a eficiência de 0,293 para
0,331 negativo. Os dois coeficientes, lado a lado, são
a evidência direta de que a identificação causal repousa sobre os
conversores (sobretudo Sorocaba e Pérola Byington, os dois com janela
posterior utilizável), e não sobre uma diferença ampla entre as redes:
a mortalidade menor é uma história dos hospitais que mudaram de
gestão, não uma diferença entre as redes OSS e Direta. E essa história
tem explicação nomeada, não apenas a constatação de que o efeito
desaparece sem os conversores: como a Subseção~\ref{sub:oqueforam}
documenta, a conversão de Sorocaba foi decidida em resposta a uma crise
de gestão, o que deprime o período anterior e faz da reversão à média a
explicação mais econômica para a melhora posterior; e no Pérola Byington
a mudança de gestão coincidiu com a mudança para um prédio novo, o que
torna o tratamento confundido. As duas alternativas não são hipóteses de
gabinete: são o que de fato aconteceu nos dois hospitais que sustentam o
resultado.

\comoler{cada linha compara duas estimativas do efeito da mesma
categoria: o círculo escuro trata cada categoria isoladamente; o
quadrado verde deixa as categorias emprestarem informação umas às
outras, na proporção da quantidade de dados de cada uma. Quanto menor
o grupo, maior o recuo do quadrado em direção ao zero: o Privado, com
3 hospitais, quase não sustenta efeito próprio. É a ilustração formal
da regra deste documento de nunca interpretar essa categoria.}

\begin{figure}[H]
\centering
\includegraphics[width=.58\textwidth]{figB02_encolhimento_categorias.png}
\caption{Efeitos de categoria na mortalidade com e sem partial
pooling: com apenas 3 hospitais, o efeito aparente do Privado encolhe
em quase 90\% quando o modelo pondera a evidência pelo tamanho do
grupo.}
\label{fig:encolhimento}
\end{figure}

\comofeito{gráfico de pontos. Cada linha é uma categoria de gestão (eixo
vertical) e o eixo horizontal traz o \textbf{efeito da categoria no
logito da mortalidade} (centrado). Cada linha traz duas estimativas: o
círculo escuro trata cada categoria isoladamente e o quadrado verde
aplica \emph{partial pooling}, em português "agrupamento parcial": o
modelo trata as categorias como aparentadas e permite que uma tome
emprestada a informação das outras, tanto mais quanto menos dados ela
própria tiver. Quanto menor o grupo, mais o quadrado recua em direção
ao zero.}

\textbf{Eficiência: a ressalva vem antes do resultado.} O DEA e a SFA
medem eficiência produtiva: quanto produto (saídas, incluída a alta
complexidade) cada hospital extrai dos seus insumos. Um dos insumos é
o \textbf{valor total real} pago pelo SUS, e ele exige a ressalva de
primeira ordem desta seção: o valor do SIH é reembolso por produção
faturada, não medida direta de recursos consumidos, e os regimes de
pagamento diferem por categoria (as OSS recebem pelo contrato de
gestão, de modo que o faturamento por AIH pode divergir do recurso
efetivo; os filantrópicos convivem com subfinanciamento e subsídio
cruzado). Parte da distância entre categorias pode refletir regime de
pagamento, não gestão. Duas verificações sustentam que a ordenação
não é artefato do insumo financeiro: a \textbf{produção física por
leito}, que não usa dinheiro nenhum, já ordena as redes do mesmo jeito
(55 saídas por leito nas OSS e 51 nos Municipais, contra 36 na Direta
e 35 nos Filantrópicos); e o DEA reestimado \textbf{só com leitos como
insumo} preserva e até amplia a vantagem (medianas corrigidas de viés
de 0,65 nas OSS contra 0,42 na Direta), com a segunda etapa
confirmando a distância
menor das OSS à fronteira (coeficiente de 1,03 negativo no logaritmo
da distância, intervalo bootstrap de 1,67 a 0,50 negativos, inteiro
longe do zero). A ressalva fica assim resolvida \textbf{no nível das
categorias} para a distância entre OSS e as demais: o afastamento da
OSS na frente não depende do insumo financeiro em nenhuma das
verificações. Já a ordem interna das outras três categorias é frágil,
como os números detalhados adiante mostram, e \textbf{os dois métodos
discordam sobre ela}. No DEA, Direta e Público Municipal ficam separados
por 0,002 (medianas corrigidas de viés de 0,634 e 0,632), um empate; na
SFA, quem empata é outro par, Direta e Filantrópico, por 0,004 (0,788
contra 0,784), enquanto o \textbf{Público Municipal aparece em 0,912},
isto é, 0,125 acima da Direta. Os dois métodos põem a OSS na frente e o
Privado no fim, e \textbf{trocam entre si a segunda e a terceira
posições}. Abaixo da OSS, portanto, o que os dados sustentam é um
\textbf{empate entre Direta e Público Municipal no DEA} (que é o
instrumento em que as conclusões de eficiência deste documento se apoiam,
pelo motivo estabelecido adiante), com o
Filantrópico atrás, e não uma hierarquia firme de quatro posições. Ela
\textbf{não} se estende ao hospital individual: a correlação de ordem
entre os escores da versão financeira e da versão só com leitos, ambas
corrigidas de viés, é de 0,51 por hospital-ano (0,48 entre as medianas
por hospital), o ranking hospital a hospital embaralha quando o
dinheiro sai do modelo, e a classificação de eficiência de um hospital
específico só deve ser citada depois de conferida nas duas versões.

\newpage
Com a ressalva posta, o resultado: na eficiência corrigida de viés
(bootstrap de Simar e Wilson, 2.000 réplicas por ano), as medianas do
período são \textbf{OSS 0,733, Direta 0,634, Público Municipal 0,632 e
Filantrópico 0,595}. Direta e Público Municipal praticamente empatam
(diferença de 0,002), sugerindo três patamares em vez de quatro: OSS à
frente, Direta e Municipal próximas no meio, Filantrópico atrás. Na
regressão da distância à fronteira sobre
categoria, complexidade, porte e ano, a vantagem OSS equivale a
\textbf{cerca de 25\% menos distância à fronteira que a Direta}
(exp(-0,293) menos 1), estável retirando o biênio pandêmico (0,272
negativo), os cinco conversores (0,331 negativo) ou restringindo ao
suporte comum (0,319 negativo), as quatro estimativas com intervalo
inteiro abaixo do zero. Ao contrário da mortalidade,
\textbf{a eficiência
maior das OSS é um traço da rede}, presente na comparação pura entre
hospitais distintos. Vale relembrar aqui, porque a frase anterior é o
tipo de afirmação que viaja sozinha, que essa conclusão é firme
\textbf{no nível das categorias} e não desce ao hospital individual: a
correlação de ordem entre a versão financeira e a versão só com leitos
é de apenas 0,51, como o parágrafo anterior registrou. E o receio de os 3 hospitais Privados
distorcerem a fronteira não se confirma: retirados, os escores dos
demais praticamente não se movem (correlação de ordem 1,000), porque
o perfil deles opera longe da fronteira.

\comoler{como nos violinos das seções anteriores: a silhueta engorda
onde há mais hospitais e o número é a mediana do grupo. Os dois painéis
concordam nas pontas, com a OSS na frente e o Privado no fim, e
\textbf{discordam no meio}: no painel (a), do DEA, a Direta (0,634) e o
Município (0,632) empatam; no painel (b), da fronteira estocástica, o
Município (0,912) sobe bem acima da Direta (0,788). O painel (b) é
diagnóstico, não medida: pelo motivo estabelecido adiante, nenhum número
dele é citável. O Privado aparece só como registro.}

\begin{figure}[H]
\centering
\subfig{figE01_violino_te_categoria.png}{Eficiência DEA corrigida de
viés}
\subfig{figS01_te_categoria.png}{Eficiência técnica da fronteira
estocástica}
\caption{Eficiência por categoria nos dois métodos. O painel (a), do DEA,
é a medida em que este documento se apoia. O painel (b), da fronteira
estocástica, é \textbf{diagnóstico e não medida}: a concentração das OSS
junto ao valor 1 é uma degeneração numérica do modelo, e a fronteira que
o gerou tem a monotonicidade violada em 94,9\% das observações, de modo
que \textbf{nenhum valor do painel (b) é citável, em nenhuma categoria}
(o texto explica). Como nos violinos da
Seção~\ref{sec:categorias}, onde a silhueta é mais gorda há mais
hospitais com aquele valor de eficiência; a largura não é percentual nem
contagem, e não se compara entre categorias.}
\label{fig:eficiencia}
\end{figure}

\comofeito{gráficos de violino, um por categoria de gestão no eixo
horizontal. O eixo vertical traz a \textbf{eficiência técnica} (entre 0 e
1, quanto mais perto de 1 mais próximo da fronteira de melhores
práticas): no painel (a), a eficiência DEA corrigida de viés, e no (b), a
eficiência da fronteira estocástica (SFA). Cada violino é a distribuição
da categoria e o número é a mediana. No painel (b) a concentração das OSS
junto a 1 é uma degeneração numérica do modelo, e nenhum valor do painel
é citável, em nenhuma categoria, pelo motivo estabelecido no texto.}

\newpage
\textbf{A contraparte paramétrica, e o que ela não é.} A
fronteira estocástica (Translog (uma forma matemática mais flexível
que a alternativa clássica, Cobb Douglas, por não obrigar os insumos a
se combinarem sempre na mesma proporção), preferida ao Cobb Douglas
com folga larga: a estatística de razão de verossimilhança entre as duas
formas é de 567,0 com 6 graus de liberdade, com p da ordem de
$10^{-119}$; e 95,5\% da variância composta é atribuída à ineficiência)
põe as OSS na mesma ponta de cima em que o DEA as põe, e traz
três honestidades a registrar, que juntas retiram dela o estatuto de
medida:
as OSS ficam tão coladas à fronteira paramétrica que o coeficiente
delas na equação de ineficiência perde identificação (o análogo da
separação perfeita da Seção~\ref{sec:estrategia}), de modo que nem o
nível nem a posição das OSS na SFA são citáveis,
e a não-identificação aparece igual nas quatro amostras e nas
três sensibilidades, com o ponto oscilando entre $-772$ e $-1.971$ e o
erro-padrão entre 400 e 1.896, que é a assinatura do problema e não uma
faixa de estimativas;\footnote{Numa das amostras (sem o biênio
pandêmico) o valor sai com $p$ de $8{,}3\times10^{-7}$. Isso \emph{não}
é precisão: sob separação a verossimilhança é praticamente plana numa
região distante de zero, e o intervalo de Wald calculado ali não tem
interpretação. Um $p$ pequeno sobre um parâmetro degenerado mede o
formato da aproximação, não o dado.} \textbf{o insumo de UTI é mal
comportado na forma Translog}, com a monotonicidade violada em
\textbf{94,9\% das observações} (97,8\% entre as que têm leito de UTI
habilitado, tenham-no usado ou não no ano; 96,9\% se o corte for ter
registrado diária de UTI),
contra 99,7\% de respeito para leitos e 99,6\% para valor: herança
do logaritmo de um mais UTI com 19\% de zeros e da colinearidade alta
com leitos totais e valor, e o motivo pelo qual qualquer leitura fina
de UTI nesta fronteira precisa da sensibilidade que a remove; e, na
outra ponta, Direta (0,788) e Filantrópico (0,784)
praticamente empatam em ponto estimado, a mesma proximidade que o DEA
mostra entre Direta e Município (sem intervalo de confiança calculado
em nenhum dos dois métodos para confirmar ou afastar sobreposição). A
correlação de ordem entre os dois conjuntos de escores é de
\textbf{0,62} (0,63 excluindo as OSS, de modo que a semelhança não
depende do grupo problemático). Ela mede quanto os dois procedimentos se
parecem quando aplicados aos mesmos dados, e \textbf{não é corroboração
de um pelo outro}: os dois consomem a mesma informação, e um deles, pelo
motivo estabelecido adiante, não descreve uma função de produção.
A semelhança não cresce de forma gradual: fica achatada entre 0,50 e
0,55 de 2015 a 2019, salta para a faixa de 0,66 a 0,74 já em 2020 e se
mantém nesse patamar mais alto até 2025, sem tendência adicional, um
degrau que coincide com o início da pandemia, não uma trajetória
ascendente ao longo da década.

\textbf{A equação de ineficiência, e por que nenhum coeficiente dela é
citável.} A fronteira estocástica estima, em paralelo à produção, uma
segunda equação que explica \emph{quanta} ineficiência cada hospital
carrega. O coeficiente das OSS ali é o número degenerado já discutido; a
equação tem \textbf{outros sete}, e \textbf{cinco} deles são fortemente
significativos. Como a variável explicada é a ineficiência, sinal
negativo significa mais eficiente.

A tabela é reproduzida porque omiti-la seria pior: ela contém o único
resultado desta fase que o documento chegou a não mostrar, e resultado
que desaparece depois de encontrado é omissão. Mas ela vem com uma
advertência que vale para cada uma das suas linhas e que precisa ser lida
\emph{antes} dos números. \textbf{A monotonicidade da fronteira que gerou
estes coeficientes está violada em 94,9\% das observações}, como o
parágrafo anterior registrou. Uma fronteira violada em quase toda a
amostra \textbf{não se comporta como função de produção}, e significância
mede a precisão de uma estimativa, não a validade da estrutura de onde
ela sai: coeficiente preciso extraído de fronteira mal-comportada é
preciso sobre um objeto que não é o que o nome diz. Por isso, e não por
conveniência de resultado, \textbf{nenhum coeficiente desta tabela é
interpretável como magnitude de ineficiência, e nenhum entra em conclusão
alguma deste documento}.

A violação também não se concentra em categoria nenhuma, o que impede
tratá-la como problema de uma linha só: ela vai de 93,6\% no Filantrópico
a \textbf{99,3\% na administração Direta}, que é justamente a referência
de todos os contrastes da tabela. Cinco vírgula sete pontos de amplitude
sobre uma base de 93,6\% é defeito geral da forma funcional. A fronteira
inteira está desqualificada como medida, não uma parte dela.

\begin{table}[H]
\centering
\caption{Equação de ineficiência da fronteira estocástica (Battese e
Coelli 1995). Referência: administração Direta, porte médio.
Coeficiente negativo indica \emph{menos} ineficiência.
\textbf{Tabela tarjada: a monotonicidade da fronteira que produziu estes
coeficientes está violada em 94,9\% das observações, de 93,6\% no
Filantrópico a 99,3\% na própria Direta. Nenhum valor desta tabela é
interpretável como magnitude nem entra em conclusão deste documento; ela
é diagnóstico da inadequação da forma Translog.}}
\label{tab:equacaoinef}
\small
\begin{tabular}{lrrr}
\toprule
Termo & Coeficiente & Erro-padrão & $p$ \\
\midrule
Público Municipal        & $-5{,}454$ & 1,020 & $9{,}0\times10^{-8}$ \\
Filantrópico             & $-0{,}244$ & 0,163 & 0,134 \\
Privado                  & $+2{,}474$ & 0,396 & $4{,}2\times10^{-10}$ \\
Complexidade (padronizada) & $-2{,}159$ & 0,325 & $3{,}2\times10^{-11}$ \\
Grande Porte             & $+2{,}470$ & 0,333 & $1{,}3\times10^{-13}$ \\
Especial                 & $+7{,}401$ & 0,917 & $7{,}0\times10^{-16}$ \\
Tendência linear         & $-0{,}010$ & 0,023 & 0,668 \\
\bottomrule
\end{tabular}
\end{table}

\textbf{O que a tabela mostra, lida como diagnóstico e não como medida.}
O coeficiente do Público Municipal ($-5{,}454$, $p$ de
$9{,}0\times10^{-8}$) é o resultado que faltava ao documento, e é também
o que melhor ilustra a advertência acima. Ele é preciso, e aponta o
Público Municipal como menos ineficiente que a administração Direta.
\textbf{O DEA, que é o instrumento adequado, não confirma essa ordem, e
tampouco a refuta.} Nas medianas corrigidas de viés impressas acima a
Direta (0,634) fica à frente do Município (0,632); nas médias o Município
fica à frente (0,644 contra 0,620). Ano a ano, no escore corrigido de
viés, o Município supera a Direta em seis dos onze anos, e os seis são de
2020 a 2025. No estrato de alta complexidade da
Seção~\ref{sec:segmentacao}, que é onde a comparação entre redes é
viável, o Município fica à frente; no estrato intermediário, atrás.
\textbf{A ordem entre estas duas redes depende de qual escore e de qual
estatística se escolhe, e por isso este documento não a afirma em direção
nenhuma.} O que há é um sinal recorrente e sem corroboração por
instrumento adequado, o que é diferente tanto de achado quanto de nada.

As demais linhas têm o mesmo estatuto. O Filantrópico não se distingue da
Direta ($p$ de 0,134). O porte carrega o maior coeficiente da tabela
(Especial, $+7{,}401$, quase três vezes o do Grande Porte) enquanto a
complexidade anda no sentido oposto ($-2{,}159$), e não há tendência
temporal ($p$ de 0,668). Lidos como \emph{ordens de grandeza relativas
dentro de uma fronteira que se sabe mal especificada}, esses sinais
sugerem que a estrutura física do hospital move a ineficiência medida
mais do que a rede que o administra. Lidos como magnitudes, não dizem
nada, pelo motivo já enunciado.

\textbf{A regra, enunciada uma vez para todo o documento.} Nenhum
coeficiente e nenhum nível da fronteira estocástica é citável, em nenhuma
categoria de gestão, nem aqui nem nas seções seguintes. A SFA permanece
no documento como \textbf{diagnóstico da própria inadequação da forma
Translog a esta tecnologia de produção}, e a violação em 94,9\% é ela
mesma um achado: diz que a Translog não descreve bem o modo como estes
hospitais transformam insumo em produto. \textbf{As conclusões de
eficiência deste documento apoiam-se no DEA.}

\comoler{quanto mais a nuvem sobe da esquerda para a direita, mais os
dois procedimentos se parecem ao ordenar os hospitais. A faixa verde
colada no topo é a degeneração das OSS descrita no texto (o modelo
paramétrico as coloca todas na fronteira); mesmo descontando essa faixa,
a semelhança permanece. Semelhança não é confirmação: o eixo vertical vem
de uma fronteira cuja monotonicidade está violada em 94,9\% das
observações, e nenhum valor dele é citável.}

\begin{figure}[H]
\centering
\includegraphics[width=.58\textwidth]{figS02_dispersao_sfa_dea.png}
\caption{Semelhança entre os dois procedimentos de eficiência: cada ponto
é um hospital num ano, posicionado pelo escore DEA (horizontal) e pela
eficiência SFA (vertical). O eixo vertical é diagnóstico, não medida:
nenhum valor de SFA é citável.}
\label{fig:convergencia}
\end{figure}

\comofeito{gráfico de dispersão. Cada ponto é um hospital em um ano,
posicionado pelo \textbf{escore de eficiência DEA} no eixo horizontal e
pela \textbf{eficiência técnica SFA} no vertical (ambos entre 0 e 1).
Quanto mais a nuvem sobe da esquerda para a direita, mais os dois métodos
(de naturezas opostas) concordam sobre quem é eficiente. A faixa
verde colada no topo é a degeneração das OSS descrita no texto.}

\newpage
\subsection*{Como os modelos de ocupação estão especificados}
\label{sub:etapa5}

Os modelos de ocupação desta fase tomam três decisões de especificação que
convém declarar, porque nas três havia alternativa plausível e a escolha
muda o que o número significa.

\textbf{O limiar de aparagem} da cauda superior é o percentil 99 da
subamostra em que a ocupação de UTI é \emph{definida}, e não do painel
completo. Calculado sobre o painel completo, o limiar seria contaminado
pelas observações de hospitais sem UTI, que não são ocupação baixa: são
ausência de ocupação (Seção~\ref{sec:categorias}).

\textbf{A amostra} do modelo de UTI é definida pela capacidade instalada,
isto é, ter leito de UTI habilitado, independentemente de tê-lo usado no
ano. O critério alternativo (restringir às observações com ocupação
positiva) selecionaria a amostra pela própria variável que se quer
explicar.

\textbf{A ocupação de internação não é aparada acima de 100\%}, porque esse
excesso não é ruído de cauda: é a assinatura aritmética da superlotação
descrita na Seção~\ref{sec:segmentacao}. Apará-lo apagaria o sinal.

Definir a amostra pela capacidade mantém dentro dela as observações de
hospitais com UTI habilitada e nenhuma diária no ano. Como a família Gama
não admite valor zero, a especificação principal é o \textbf{modelo em duas
partes com Gama na parte contínua}, que acomoda essa massa em zero num
componente próprio e preserva a função de variância da Gama.

\textbf{O resultado.} As OSS operam a UTI \textbf{30,4\% mais
intensamente} que a administração Direta, com intervalo de 95\% entre
$+5{,}7\%$ e $+60{,}8\%$. A ocupação de internação aponta na mesma direção,
$+11{,}5\%$, mas com o intervalo cruzando o zero, de modo que ali não se
pode afirmar diferença.

\subsubsection*{A escolha de família move mais o número do que os defeitos}

Entre todas as decisões de modelagem, a que mais desloca a estimativa é a
\textbf{família de distribuição}, mais do que qualquer das três decisões
de especificação acima. Por isso as cinco especificações consideradas são
reportadas juntas.

\begin{table}[H]
\centering
\caption{Intensidade de uso da UTI, OSS contra administração Direta, sob
cinco especificações. Efeito e intervalo em percentual.}
\label{tab:familiasuti}
\small
\begin{tabular}{lrrrl}
\toprule
Especificação & $n$ & Efeito & IC 95\% & Zero \\
\midrule
\textbf{Duas partes, Gama (principal)} & 2.489 & $\mathbf{+30{,}4\%}$ & $[+5{,}7\%;\ +60{,}8\%]$ & exclui \\
Gama, apenas os positivos              & 2.482 & $+30{,}4\%$ & $[+5{,}7\%;\ +60{,}8\%]$ & exclui \\
Tweedie sem os zeros                   & 2.482 & $+27{,}0\%$ & $[+7{,}6\%;\ +49{,}8\%]$ & exclui \\
Tweedie com os zeros                   & 2.489 & $+26{,}7\%$ & $[+6{,}7\%;\ +50{,}4\%]$ & exclui \\
\textbf{Lognormal sem os zeros}        & 2.482 & $+25{,}9\%$ & $[-0{,}3\%;\ +59{,}0\%]$ & \textbf{cruza} \\
\bottomrule
\end{tabular}
\end{table}

A amplitude entre as cinco é de \textbf{4,5 pontos percentuais}, menor que a
largura do intervalo de confiança de qualquer uma delas. A direção é a mesma
nas cinco, e \textbf{quatro excluem o zero}.

\textbf{A quinta é a lognormal, e ela precisa ser nomeada}, porque
``significativo em quatro de cinco'' sem dizer qual é a exceção não é
robustez reportada. \textbf{E ela não discorda.} O ponto dela está a
\textbf{0,0087} em pontos de logaritmo do Tweedie estimado sobre exatamente
a mesma amostra, distância menor que a que separa quaisquer duas das
outras famílias. O que muda é a \textbf{precisão}: o erro-padrão da
lognormal é 41\% maior que o do Tweedie e 11\% maior que o da Gama.
\textbf{Ela perde a significância por variância, não por deslocamento}, e
por margem mínima: o limite inferior fica a três milésimos do zero.

A razão é que a lognormal \textbf{estima outra quantidade}. As demais
modelam a média condicional; a lognormal modela a mediana condicional. Sob
assimetria à direita, que é a forma desta distribuição, o efeito sobre a
mediana é sistematicamente mais atenuado e menos preciso que o efeito sobre
a média. É \emph{estimar diferente}, não \emph{estimativa discordante}, e
isso fortalece a leitura de robustez em vez de qualificá-la: a única
especificação que não atinge significância é também a única que não está
medindo a mesma coisa que as outras, e ainda assim entrega ponto
equivalente.

\comofeito{um alerta sobre a coluna de ajuste que não aparece nesta tabela:
o critério de informação da lognormal (2.135) é numericamente muito menor
que o das demais (entre 21.608 e 22.695), e \textbf{isso não significa
ajuste melhor}. O da lognormal é calculado sobre o logaritmo da resposta e o
das outras sobre a resposta, escalas diferentes. O critério é comparável
\emph{dentro} de cada escala, nunca entre elas.}

\newpage
\textbf{O que fica para a leitura do documento.} A segunda fase não
altera nenhuma conclusão; ela redistribui os pesos. O \textbf{regime
operacional mais intenso e a eficiência produtiva maior das OSS são
traços robustos da rede}, mas é preciso dizer \emph{qual} componente do
regime operacional resiste a quê. \textbf{Produção e intensidade de uso
da UTI excluem o zero nas quatro amostras}, inclusive sem os
conversores (produção de $+0{,}26$ a $+0{,}49$ em pontos de log;
intensidade de UTI de $+0{,}19$ a $+0{,}50$). \textbf{A ocupação de
internação, não:} o intervalo cruza o zero em duas das quatro, na
especificação principal ($-0{,}041$ a $+0{,}277$) e sem o biênio
pandêmico ($-0{,}077$ a $+0{,}260$), e só o exclui sem os
conversores e no suporte comum. Escrever que o regime operacional
inteiro é robusto ``em todas as amostras'' seria estender à ocupação
uma robustez que ela não tem, e que as duas passagens anteriores desta
mesma Seção já haviam registrado; o \textbf{TMP menor} reaparece
com a mesma magnitude na especificação hierárquica, o que confirma que
o resultado não dependia da forma do modelo, mas não altera a
limitação que vem dos cinco conversores, e ele segue sem corroboração
nos testes de poder limitado (permutação p igual a 0,153, bootstrap
selvagem 0,324, controle sintético degenerado); o \textbf{faturamento} segue sem
diferença em qualquer método; e a \textbf{mortalidade menor}
concentra-se inteiramente na trajetória dos cinco conversores, o
que a aproxima das ressalvas causais das
Seções~\ref{sec:efeitogestao} e~\ref{sec:inferenciafinal} e a afasta
de uma leitura de superioridade da rede OSS nesse desfecho.
